\begin{abstract}
Governance frameworks ask AI providers and auditors for
\emph{documented evaluation evidence}, and perturbation-based
construct-validity audits are a common form of that evidence. We
argue the audits are themselves fragile: their conclusions can be
silently manufactured by implementation details that readers
cannot see in the reported numbers. We name five classes of
pipeline failure and demonstrate each in a self-audit over safety
benchmarks and open-weight instruction-tuned models. Under a
unified six-point due-diligence gate, every cell lands in a
non-confirmatory bucket, and no cell reaches \emph{confirmatory}.
The evidence here is a single two-model, five-benchmark case
study, and F1--F5 is an illustrative, deliberately non-exhaustive
starting taxonomy---not a comprehensive partition of audit
failures. We position the gate as a withholding and disclosure
protocol for assurance-grade evidence, supplementary to (not a
replacement for) classical construct-validity evidence, and not as
a route to benchmark-validity verdicts.
\end{abstract}
